
\exercise{Traffic: Simple centralities}{3}
Consider the following road network
\picm{Gritting}{}
\subquestion Compute the degree centrality of all nodes.

\solution
This one is straight forward. 
\begin{center}
A: 3, B: 4, C: 4, D: 3, E: 4, F: 4, G: 4, H: 3, X: 1
\end{center}
The result is reasonable but not very useful. It tells us that crossroads are more important than dead ends, but does not help us distinguish the relative importance of the degree 4 nodes. 

%%%%%%%%%

\subquestion
Compute the closeness centrality of all nodes. (All links have length 1)

\solution 
For this one we need to find the shortest paths in the network. For this small example visual inspection is easier than Dijkstra. We find 
\begin{center}
\begin{tabular}{ c c c | c c c }
Trip & Distance & Intermediate nodes & Trip & Distance & Intermediate nodes \\\hline
A-B & 1 & - &
A-C & 1 & - \\
A-D & 2 & B &
A-E & 2 & C \\
A-F & 3 & B,G / C,E &
A-G & 2 & B \\
A-H & 3 & B,D / C,E &
A-X & 1 & - \\
B-C & 1 & - &
B-D & 1 & - \\
B-E & 2 & C &
B-F & 2 & G \\
B-G & 1 & - &
B-H & 2 & D \\
B-X & 2 & A &
C-D & 2 & B \\
C-E & 1 & - &
C-F & 2 & E \\
C-G & 2 & B &
C-H & 2 & E \\
C-X & 2 & A &
D-E & 2 & H \\
D-F & 2 & H / G &
D-G & 1 & - \\
D-H & 1 & - &
D-X & 3 & A,B \\
E-F & 1 & - &
E-G & 2 & F \\
E-H & 1 & - &
E-X & 3 & C,A \\
F-G & 1 & - &
F-H & 1 & - \\
F-X & 4 & E,C,A / G,B,A &
G-H & 2 & F / D \\
G-X & 3 & A,B &
H-X & 4 & A,B,D / A,C,E
\end{tabular}
\end{center}
We can now compute the centralities
\eqa{
c_{\rm A}&=&\frac{8}{1+1+2+2+3+2+3+1}=\frac{8}{15}\approx 0.53 \\
c_{\rm B}&=&\frac{8}{1+1+1+2+2+1+2+2}=\frac{8}{12}= 0.75 \\
c_{\rm C}&=&\frac{8}{1+1+2+1+2+2+2+2}=\frac{8}{13}\approx 0.61 \\
c_{\rm D}&=&\frac{8}{2+1+2+2+2+1+1+3}=\frac{8}{13}\approx 0.61 \\
c_{\rm E}&=&\frac{8}{2+2+1+2+1+2+1+3}=\frac{8}{14}\approx 0.57\\
c_{\rm F}&=&\frac{8}{3+2+2+2+1+1+1+4}=\frac{8}{16}= 0.50\\
c_{\rm G}&=&\frac{8}{2+1+2+1+2+1+2+3}=\frac{8}{14}\approx 0.57\\
c_{\rm H}&=&\frac{8}{3+2+2+1+1+1+2+4}=\frac{8}{16}= 0.50\\
c_{\rm X}&=&\frac{8}{1+2+2+3+3+4+3+4}=\frac{8}{22}\approx 0.36 
}
This is soo much more work than the spectral centrality! Note that again X lands in last place which is reasonable, while B is the most important node. 

%%%%%%%

\subquestion
Compute the betweenness centrality of all nodes.

\solution
Thankfully, we can reuse the table from the previous question. We just need to count occurrences of intermediate nodes. The challenge is to not let the bookkeeping get out of hand. For this purpose I made another table

\begin{center}
\begin{tabular}{c c c c}
Node & On Paths (no alternative) & On Path (one alternative) & Centrality \\\hline
A & B-X,C-X,D-X,E-X,F-X,G-X,H-X & - & 7 \\
B & A-D,A-G,C-D,C-G,D-X,G-X & A-F,A-H,F-X,H-X & 8 \\
C & A-E,B-E,E-X & A-F,A-H,F-X,H-X & 5 \\
D & B-H & A-H,F-X,G-H,H-X & 3 \\
E & C-F,C-H & A-F,A-H,F-X,H-X & 4 \\
F & E-G & G-H & 1.5 \\
G & B-F & A-F,D-F & 2   \\
H & D-E & D-F & 1.5 \\
X & - & - & 0 \\  
\end{tabular}
\end{center}
Again we end up with B as the most important node and X as the least important which makes sense. A surprise is maybe the relatively high value for A. This highlights A as a bottleneck. Generally, betweenness centrality is used when we are interested in finding such bottlenecks. 

